% Bloom levels: remember, understand, apply, analyze, evaluate, create.
\begin{problema}\label{prob:ddb67d0}
A battery manufacturer claims that the mean lifetime of its product is $\mu_0=500$ hours. An independent laboratory suspects that the true mean lifetime is \emph{lower}. Formulate the null hypothesis $H_0$ and the alternative hypothesis $H_a$, and identify whether the test is left-tailed, right-tailed, or two-tailed.
\end{problema}

\begin{problema}\label{prob:78013c6}
A screening test for a rare disease is designed with $H_0$: ``the patient is healthy'' against $H_a$: ``the patient is sick.'' Describe in words what Type I error and Type II error represent in this specific medical context, and reason which of the two errors has more serious practical consequences.
\end{problema}

\begin{problema}\label{prob:74dc84f}
A restaurant chain claims that the average waiting time is $\mu_0=12$ minutes. A sample of $n=49$ customers gives $\bar x=12.9$ minutes, with known population standard deviation $\sigma=2.8$ minutes. Formulate the hypotheses for a two-tailed test, compute the $Z$ statistic, obtain the corresponding $p$-value, and decide at level $\alpha=0.05$.
\end{problema}

\begin{problema}\label{prob:fad6841}
For the test $H_0:\mu=\mu_0$ against $H_a:\mu>\mu_0$ with known $\sigma$ and fixed sample size $n$, derive the general expression for the power function $\operatorname{Power}(\mu_a)=1-\beta(\mu_a)$ as a function of the true value $\mu_a$, expressing it in terms of $\Phi(\cdot)$. Analyze the behavior of $\operatorname{Power}(\mu_a)$ as $\mu_a\to\mu_0$ and as $\mu_a\to\infty$.
\end{problema}

\begin{problema}\label{prob:66368eb}
A data scientist runs $m=20$ statistically independent hypothesis tests on different true null hypotheses, each at individual level $\alpha=0.05$, without applying any correction. Evaluate this practice by computing the resulting family-wise error rate (FWER), and propose and justify the corrected individual significance level $\alpha^*$ required by the Bonferroni correction to keep the global FWER at $0.05$.
\end{problema}

\begin{problema}\label{prob:d04b6da}
Design an original hypothesis-test example, with context, $H_0$, $H_a$, and one more costly type of error, different from the battery, restaurant, or medical-test examples in this section. Formulate both hypotheses, identify the direction of the test, and explain which error, Type I or Type II, would be more serious in your context and why.
\end{problema}

\begin{solproblema}[prob:ddb67d0]
The manufacturer's claimed value is the null hypothesis:
\[
	H_0:\mu=500\text{ hours},\qquad H_a:\mu<500\text{ hours}.
\]
Since the alternative points toward smaller values, this is a left-tailed test.
\end{solproblema}

\begin{solproblema}[prob:78013c6]
A Type I error means rejecting a true $H_0$: the test says the patient is sick when the patient is actually healthy, a false positive. A Type II error means not rejecting a false $H_0$: the test says the patient is healthy when the patient is actually sick, a false negative. In a disease-screening context, Type II error is usually more serious because it can delay treatment for a sick patient, whereas a false positive often leads to additional confirmatory testing.
\end{solproblema}

\begin{solproblema}[prob:74dc84f]
The hypotheses are
\[
	H_0:\mu=12,\qquad H_a:\mu\ne12.
\]
The statistic is
\[
	Z=\frac{12.9-12}{2.8/\sqrt{49}}=\frac{0.9}{0.4}=2.25.
\]
The two-tailed $p$-value is
\[
	p=2(1-\Phi(2.25))\approx2(0.01222)=0.02444.
\]
Since $0.02444<0.05$, reject $H_0$: there is evidence that the average waiting time differs from $12$ minutes.
\end{solproblema}

\begin{solproblema}[prob:fad6841]
The rejection region has the form $\bar X>c$, where
\[
	c=\mu_0+Z_\alpha\frac{\sigma}{\sqrt n}
\]
under $H_0$. For a true value $\mu_a$,
\[
	\operatorname{Power}(\mu_a)
	=P(\bar X>c\mid\mu=\mu_a)
	=1-\Phi\left(\frac{c-\mu_a}{\sigma/\sqrt n}\right)
	=1-\Phi\left(Z_\alpha+\frac{\mu_0-\mu_a}{\sigma/\sqrt n}\right).
\]
As $\mu_a\to\mu_0$, the power tends to $1-\Phi(Z_\alpha)=\alpha$. As $\mu_a\to\infty$, the argument tends to $-\infty$, so the power tends to $1$.
\end{solproblema}

\begin{solproblema}[prob:66368eb]
Under independence, the probability of no Type I error across all tests is $(1-\alpha)^m$. Thus
\[
	\operatorname{FWER}=1-(1-0.05)^{20}=1-0.95^{20}\approx0.6415.
\]
So the chance of at least one false positive is about $64.15\%$, making the uncorrected practice inappropriate. Bonferroni uses
\[
	\alpha^*=\frac{0.05}{20}=0.0025,
\]
which guarantees $\operatorname{FWER}\le20(0.0025)=0.05$.
\end{solproblema}

\begin{solproblema}[prob:d04b6da]
Example: a credit-card fraud system uses $H_0$: ``the transaction is legitimate'' against $H_a$: ``the transaction is fraudulent.'' The test is oriented toward detecting a rare event. A Type I error blocks a legitimate transaction, annoying the customer; a Type II error misses a fraudulent transaction, producing a direct financial loss. In this context, Type II error is often more costly, although too many Type I errors can also damage customer experience. The system should balance both costs explicitly.
\end{solproblema}
